
est
commutatif. Soit @@M0@@~; alors~: @@M1@@~; d'où~:
@@M2@@

Il reste donc à montrer que le morphisme de foncteurs ainsi construit est compatible aux structures de groupes filtrés.

Supposons que @@M3@@. Dire que @@M4@@ est de filtration @@M5@@ équivaut à dire que @@M6@@, se factorise par l'endomorphisme @@M7@@ de @@M8@@ défini par~: @@M9@@ si @@M10@@, et @@M11@@ si @@M12@@.
Pour montrer que @@M13@@ est compatible aux filtrations, il suffit donc de montrer que @@M14@@. Or par construction~:
@@M15@@
d'où~:
@@M16@@

Pour vérifier l'additivité de @@M17@@, il suffit de le faire pour les éléments de degré fini de @@M18@@, et de passer ensuite à la limite. Soient donc @@M19@@ deux éléments de degré @@M20@@ de @@M21@@, et @@M22@@ leur produit. On considère le diagramme commutatif~:
@@M23@@
où les flèches sont définies comme suit~: @@M24@@ opère sur @@M25@@ 
